
\thispagestyle{empty}
\begin{textblock*}{\textwidth}(-2cm,-2.5cm)
\includegraphics[height=0.1\paperheight]{logos/Bicocca_transparent.png}
\end{textblock*}
\begin{textblock*}{\textwidth}(1.9cm,-2.6cm)
\begin{flushleft}
\Large{DOCTORAL SCHOOL}\\[2mm]
\Large{UNIVERSITY OF MILANO-BICOCCA}
\end{flushleft}
\end{textblock*}
\vspace{1cm}
\centerline{\Large{Department of \textbf{Informatics, Systems and Communication}}}
\vspace{0.5cm}
\centerline{\Large{Ph. D. program in \textbf{Computer Science}, \textbf{XXXVIII}} cycle}
\vspace{2mm}
% \centerline{\Large{Curriculum of \alert{\textbf{your curriculum (if applicable)}}}}
\vspace{1cm}
% \centerline{\huge{Entity-Oriented Knowledge Discovery in Knowledge-Intensive Domains}}
% \huge{Entity-Oriented Strategies for Information Extraction and Access in the Legal Domain}
\begin{center}
  \huge{Entity-Oriented Strategies for Information Extraction and Access in Knowledge-Intensive Domains}
\end{center}

\vspace{2cm}
\centerline{\LARGE{\textbf{Riccardo Pozzi}}}
\vspace{2mm}
\centerline{\Large{Student Number 807857}}
\vspace{2cm}
\begin{flushleft}
\Large{Supervisor: \textbf{Prof. Matteo Palmonari}}\\[1.5mm]
\Large{Ph. D. Tutor: \textbf{Prof. Claudio Zandron}}\\[1.5mm]
\Large{Ph. D. Coordinator: \textbf{Prof. Leonardo Mariani}}\\[2cm]
% \LARGE{Coordinator: \alert{\textbf{Leonardo Mariani and Gianluca Della Vedova}}}
\end{flushleft}
% \vfill

\begin{flushright}
\LARGE{Academic Year \textbf{2024--2025}}
\end{flushright}
%%%%%%%%%%%%%%%%%
%%% ABSTRACT (Techincal)%%%%
\clearpage
\begin{abstract}
%    A one-page \say{technical} abstract that summarizes
%    your research work

% In certain domains, such as ..., users face difficulties ...ing with legacy technologies.
% In several cases, semantic knowledge can simplify their exploration,... activities.
% In my research, I worked first on knowledge extraction...
% Then on knowledge access and integration...


% %%%%%
% In \textbf{knowledge-intensive domains}, such as the legal and investigative sectors, users face difficulties interacting with high volumes of data. Indeed, documents such as legal judgements or instant-messaging data seized from suspects can be time-consuming to explore.
% These documents contain several mentions of \textbf{named entities}, such as people, locations, or organizations, making their extraction and organization (for example in a knowledge graph) particularly useful to simplify the exploration and interpretation of these documents by the end users (judges, lawyers, or investigators).

% My research has been directed towards \textbf{knowledge extraction and knowledge access} of the information present in domain-specific documents. While extracting and accessing knowledge imply different technical tasks, these topics are equally critical to simplify the \textbf{knowledge discovery process by end users}, that is the main topic of my research.

% With respect to \textbf{knowledge extraction}, I focused on \textbf{named entity recognition and linking}, addressing the challenges posed by \textbf{NIL entities}, those not present in existing knowledge bases. These entities are particularly important for domain-specific scenarios and have been often ignored in entity linking studies.
% In my co-authored papers we tackled this problem by publishing NIL-aware datasets, one for generain-domain textual documents and tabular data, developing baseline techniques for identifying and clustering NIL mentions.

% Subsequently, I focused on extracting and organizing entities from \textbf{domain-specific documents}, considering both the law and investigative domains, \textbf{manually labeling} them with mentions of entities and highlighting the \textbf{challenges} due to the peculiarities of the domain-specific language.
% Additionally, I studied \textbf{adaptation methods} for named entity recognition on Italian judgements, concluding that task.specific fine-tuning is a good compromise between performance and computational costs.

% Regarding \textbf{knowledge access and its integration with user interfaces}, I designed infrastructures that support entity-centric exploration, including faceted search and question answering, using knowledge graphs for representing the network of relations in the investigative domain and including \textbf{human-in-the-loop} functionalities that allow users to verify and correct machine errors.

% I also developed \textbf{ReFactX}, a scalable and efficient method for answering questions using external knowledge bases through constrained generation. This work supports users in accessing structured knowledge with minimal computational resources, especially in privacy-sensitive contexts where reliance on public APIs is not feasible.
% %%%%%

Knowledge-intensive domains such as law require accessing, integrating, and reasoning over large collections of heterogeneous documents, while meeting strict requirements on privacy, traceability, and regulatory compliance. Despite recent advances in large language models, their direct use in these settings is limited by hallucination, lack of grounding, and the legal constraints that complicate the transfer of sensitive data to external \acsp{api}. This thesis investigates how to satisfy information needs in the legal domain, including precedent retrieval, investigative search on seized data, document navigation, question answering, and statistical monitoring, under these constraints.

The work pursues three objectives. First, it quantifies to what extent general-domain \acl{ee} pipelines (\acl{ner}, \acl{nel}, NIL prediction) can be applied to legal judgments and investigative chat logs. The results show that incremental \acl{ee}, where novel entities are identified and added to the \acl{kb}, suffer from error propagation, and that NIL prediction is a major bottleneck, supporting the need for architectures that tolerate imperfect extraction. Second, it designs an entity-centric \acl{dia} that integrates heterogeneous legal sources (judgments, investigative chats, attachments) around entities, supports traceability and human oversight via error correction functionalities, remains useful despite extraction errors, and enables the use cases above. Third, it develops \refactx, a constrained-generation approach to \acl{qa} that injects facts from large knowledge bases into a \acl{llm} without retrievers or external calls, producing answers that are verifiable against grounded evidence while adding only negligible latency, thus remaining efficient and suitable for local deployment.

Together, the contributions form an integrated approach for information access in knowledge-intensive legal settings. The entity-centric \acl{dia} integrates the knowledge received from extraction services and can be paired with \aclp{ui} or \refactx to support downstream use cases, while preserving traceability, verifiability, and error-correction capability in line with the GDPR and AI Act requirements on data control and human oversight.



\end{abstract}
%%% PUBLICATIONS %%%
\pagenumbering{roman}

% Name & headline

%\switchcolumn
% Font: \fontname\font

\clearpage
\section*{List of publications}
\begin{itemize}
  \item \fullcite{ijckg2022}
  \item \fullcite{mammotab}
  \item \fullcite{BELLANDI2024}
  \item \fullcite{aixia_2023}
  \item \fullcite{pozzi_wise_2025}
  \item \fullcite{Agazzi2025DAVE}
  \item \fullcite{Pozzi2025ReFactX}
\end{itemize}

\section*{Contributions to conferences}
As a \textbf{speaker} I presented:
\begin{itemize}
	\item Ref.~\cite{ijckg2022} at the 11th International Joint Conference on Knowledge Graphs (IJCKG 2022).
	\item Ref.~\cite{aixia_2023} at the 22nd International Conference of the Italian Association for Artificial Intelligence (AIxIA 2023).
	\item Ref.~\cite{pozzi_wise_2025} at the 25th International Conference on Web Information Systems Engineering (WISE 2024).
	\item Ref.~\cite{Pozzi2025ReFactX} at the 24th International Semantic Web Conference (ISWC 2025).
\end{itemize}

As a \textbf{poster presenter}:
\begin{itemize}
    \item Poster based on Ref.~\cite{Pozzi2025ReFactX} at the Lectures on Computational Linguistics 2025, organized by the Italian Association of Computational Linguistics (AILC).
\end{itemize}

As an \textbf{author}:
\begin{itemize}
	\item Ref.~\cite{ijckg2022} in the Proceedings of the 11th International Joint Conference on Knowledge Graphs (IJCKG 2022).
	\item Ref.~\cite{mammotab} in the Proceedings of the Fourth edition of the Semantic Web Challenge on Tabular Data to Knowledge Graph Matching (SemTab 2022).
	\item Ref.~\cite{aixia_2023} in the Proceedings of the 22nd International Conference of the Italian Association for Artificial Intelligence (AIxIA 2023).
	\item Ref.~\cite{pozzi_wise_2025} in the Proceedings of the 25th International Conference on Web Information Systems Engineering (WISE 2024).
	\item Ref.~\cite{Agazzi2025DAVE} in the Proceedings of the 34th International Joint Conference on Artificial Intelligence (IJCAI 2025), Demo Track.
	\item Ref.~\cite{Pozzi2025ReFactX} in the Proceedings of the 24th International Semantic Web Conference (ISWC 2025).
\end{itemize}

%%%%%%%
\clearpage
\tableofcontents

\clearpage
\listoffigures
\clearpage
\listoftables
\clearpage
\printacronyms

\pagestyle{mystyle}

\pagenumbering{arabic}

% \addcontentsline{toc}{chapter}{Introduction}